\documentclass{article}
\usepackage[spanish]{babel}
\usepackage[T1]{fontenc}
\usepackage[utf8]{inputenc}
\usepackage{graphicx}
\usepackage{parskip}
\usepackage{listings}
\usepackage{xcolor}
\usepackage{hyperref}
 
% ── Estilo para código Java ──────────────────────────────────────────────────
\lstdefinestyle{javaStyle}{
  language=Java,
  basicstyle=\ttfamily\small,
  keywordstyle=\color{blue}\bfseries,
  stringstyle=\color{orange},
  commentstyle=\color{gray}\itshape,
  numbers=left,
  numberstyle=\tiny\color{gray},
  stepnumber=1,
  numbersep=6pt,
  backgroundcolor=\color{gray!8},
  frame=single,
  rulecolor=\color{gray!40},
  breaklines=true,
  breakatwhitespace=false,
  tabsize=4,
  captionpos=b,
  showstringspaces=false
}
 
% ── Estilo para código Dart / Flutter ────────────────────────────────────────
\lstdefinestyle{dartStyle}{
  language=Java,          % Dart no es nativo, Java da resaltado aceptable
  basicstyle=\ttfamily\small,
  keywordstyle=\color{teal}\bfseries,
  stringstyle=\color{orange},
  commentstyle=\color{gray}\itshape,
  morekeywords={class, final, const, required, async, await, Future, String,
                bool, int, List, Map, Widget, override, get, null, void,
                static, factory, extends, implements, new, return, if,
                else, true, false, late, var, dynamic, super, this,
                import, export, package},
  numbers=left,
  numberstyle=\tiny\color{gray},
  stepnumber=1,
  numbersep=6pt,
  backgroundcolor=\color{teal!6},
  frame=single,
  rulecolor=\color{teal!30},
  breaklines=true,
  breakatwhitespace=false,
  tabsize=2,
  captionpos=b,
  showstringspaces=false
}
 
\title{Patrones de Dise\~no en mi proyecto:\\Simulador de Combates de Dragon Ball}
\author{Diego}
\date{Mayo 5, 2026}
 
\begin{document}
 
\maketitle
\tableofcontents
\newpage
 
% ============================================================================
\section{Introducci\'on}
% ============================================================================
 
Los patrones de dise\~no son soluciones probadas a problemas comunes en el
desarrollo de software. Se dividen en tres grupos principales: los
\textbf{creacionales}, que tienen que ver con c\'omo se crean los objetos;
los \textbf{estructurales}, que hablan de c\'omo se organizan entre s\'i, y
los \textbf{de comportamiento}, que definen c\'omo interact\'uan y se comunican.
 
En este documento explico los patrones que utilic\'e en mi proyecto,
\textit{Simulador de peleas de Dragon Ball}, el cual est\'a compuesto por:
 
\begin{itemize}
  \item \textbf{Backend:} API REST en Java con Spring Boot 3, Spring Security y
        JWT, conectada a una base de datos PostgreSQL en Supabase y desplegada
        en Render.
  \item \textbf{Frontend:} Aplicaci\'on m\'ovil en Flutter (Dart) que consume la
        API y sube im\'agenes directamente a Cloudinary.
\end{itemize}
 
% ============================================================================
\section{Patrones Creacionales}
% ============================================================================
 
Estos patrones se encargan de controlar c\'omo se crean los objetos dentro de
una aplicaci\'on. Puede parecer que eso es lo m\'as simple del mundo (solo haces
\texttt{new Objeto()} y listo), pero en proyectos m\'as grandes eso puede
volverse un problema.
 
% ----------------------------------------------------------------------------
\subsection{Singleton}
% ----------------------------------------------------------------------------
 
La idea es que de una clase solo exista una \'unica instancia en todo el
programa. Imag\'inate que tienes una clase que maneja la conexi\'on a la base
de datos: no tiene sentido crear cien conexiones distintas, con una sola
alcanza.
 
En el proyecto utilic\'e el Singleton para el \texttt{AuthService} en Flutter,
donde guardo la sesi\'on del usuario. No quer\'ia que hubiera varios objetos
manejando el token al mismo tiempo porque eso generar\'ia inconsistencias.
 
\subsubsection*{Implementaci\'on en el Frontend (Dart / Flutter)}
 
La clase \texttt{AuthService} declara un campo est\'atico \texttt{instance}
que es la \'unica instancia posible. El constructor privado \texttt{AuthService.\_()}
impide que cualquier otra parte del c\'odigo pueda crear un objeto nuevo.
 
\begin{lstlisting}[style=dartStyle, caption={AuthService -- patr\'on Singleton (lib/services/auth\_service.dart)}]
class AuthService {
  static const _keyToken   = 'auth_token';
  static const _keyUsuario = 'auth_usuario';
 
  // Unica instancia global; se accede con AuthService.instance
  static final AuthService instance = AuthService._();
  AuthService._();   // Constructor privado: nadie mas puede instanciar
 
  Usuario? _usuario;
 
  Usuario? get usuario       => _usuario;
  bool     get estaAutenticado => _usuario != null;
  String?  get token           => _usuario?.token;
 
  // Cargar la sesion guardada al arrancar la app
  Future<void> cargarSesion() async {
    final prefs = await SharedPreferences.getInstance();
    final json  = prefs.getString(_keyUsuario);
    if (json != null) {
      try {
        _usuario = Usuario.fromJson(jsonDecode(json));
      } catch (_) {
        await cerrarSesion();
      }
    }
  }
  // ... resto de metodos (login, registrar, cerrarSesion)
}
\end{lstlisting}
 
En \texttt{main.dart} se accede al singleton antes de arrancar la app:
 
\begin{lstlisting}[style=dartStyle, caption={Uso del Singleton en main.dart}]
void main() async {
  WidgetsFlutterBinding.ensureInitialized();
  // Se usa la instancia unica para cargar el token guardado
  await AuthService.instance.cargarSesion();
  runApp(const DragonBallApp());
}
\end{lstlisting}
 
Cualquier widget del proyecto puede obtener el token actual con
\texttt{AuthService.instance.token} sin necesidad de pasar ning\'un objeto
entre pantallas.
 
La desventaja que le veo es que a veces complica las pruebas, porque el
estado global puede interferir entre tests, pero en general es muy \'util.
 
% ----------------------------------------------------------------------------
\subsection{Factory Method}
% ----------------------------------------------------------------------------
 
Este patr\'on b\'asicamente dice: en lugar de crear un objeto directamente,
delegas esa responsabilidad a un m\'etodo o clase especial que decide qu\'e
tipo de objeto crear.
 
Lo us\'e de forma informal cuando el backend decide qu\'e tipo de respuesta
construir dependiendo del resultado de la operaci\'on. Adem\'as, en Dart los
constructores nombrados \texttt{fromJson} son una variante directa del
Factory Method: delegan la creaci\'on del objeto a un m\'etodo que sabe c\'omo
interpretar los datos entrantes.
 
\subsubsection*{Implementaci\'on en el Frontend (Dart / Flutter)}
 
Los modelos \texttt{Combate} y \texttt{Usuario} usan constructores
\texttt{factory} que act\'uan como f\'abrica: reciben un \texttt{Map} JSON y
devuelven una instancia correctamente inicializada.
 
\begin{lstlisting}[style=dartStyle, caption={Factory Method en el modelo Combate (lib/models/combate.dart)}]
class Combate {
  final int?   id;
  final int    anio;
  final String intensidad;
  final String tipo;
  final String lugar;
  final String creador;
  final int?   usuarioId;
  final String? resultado;
  final String? fechaCreacion;
  final ImagenCombate? imagen;
 
  Combate({ this.id, required this.anio, required this.intensidad,
            required this.tipo, required this.lugar, required this.creador,
            this.usuarioId, this.resultado, this.fechaCreacion, this.imagen });
 
  // Factory Method: decide como construir el objeto segun el JSON recibido
  factory Combate.fromJson(Map<String, dynamic> j) => Combate(
    id:           j['id'],
    anio:         j['anio'],
    intensidad:   j['intensidad'],
    tipo:         j['tipo'],
    lugar:        j['lugar'],
    creador:      j['creador'],
    usuarioId:    j['usuarioId'],
    resultado:    j['resultado'],
    fechaCreacion: j['fechaCreacion'],
    imagen: j['imagen'] != null
        ? ImagenCombate.fromJson(j['imagen'])
        : null,   // <-- decision segun el dato disponible
  );
}
\end{lstlisting}
 
\subsubsection*{Implementaci\'on en el Backend (Java / Spring Boot)}
 
En el backend, \texttt{AuthService} construye distintos tipos de respuesta
seg\'un la operaci\'on (registro o login), encapsulando la l\'ogica de creaci\'on
del objeto \texttt{AuthDtos.AuthResponse}.
 
\begin{lstlisting}[style=javaStyle, caption={Creaci\'on de respuesta en AuthService (service/AuthService.java)}]
@Service
public class AuthService {
 
    private final UsuarioRepository repo;
    private final PasswordEncoder   encoder;
    private final JwtUtil           jwtUtil;
 
    public AuthDtos.AuthResponse registrar(AuthDtos.RegisterRequest req) {
        if (repo.findByEmail(req.getEmail()).isPresent())
            throw new IllegalArgumentException("Email ya registrado");
 
        Usuario u = new Usuario();
        u.setNombre(req.getNombre());
        u.setEmail(req.getEmail());
        u.setPassword(encoder.encode(req.getPassword()));
        repo.save(u);
 
        String token = jwtUtil.generarToken(u.getEmail());
        // Factory: construye la respuesta segun el contexto (registro)
        return new AuthDtos.AuthResponse(
                token, u.getId(), u.getNombre(), u.getEmail());
    }
 
    public AuthDtos.AuthResponse login(AuthDtos.LoginRequest req) {
        Usuario u = repo.findByEmail(req.getEmail())
                .orElseThrow(() -> new IllegalArgumentException(
                        "Usuario no encontrado"));
        if (!encoder.matches(req.getPassword(), u.getPassword()))
            throw new IllegalArgumentException("Contrasena incorrecta");
 
        String token = jwtUtil.generarToken(u.getEmail());
        // Factory: construye la respuesta segun el contexto (login)
        return new AuthDtos.AuthResponse(
                token, u.getId(), u.getNombre(), u.getEmail());
    }
}
\end{lstlisting}
 
% ----------------------------------------------------------------------------
\subsection{Builder}
% ----------------------------------------------------------------------------
 
El Builder sirve cuando tienes objetos que se construyen paso a paso y
que tienen muchos campos opcionales. En vez de tener un constructor con
veinte par\'ametros (que nadie entiende en qu\'e orden van), usas un objeto
intermedio que vas configurando.
 
En mi proyecto, cuando se construye un objeto \texttt{Combate} con lugar,
a\~no, intensidad, tipo, imagen y resultado, se podr\'ia usar un Builder para
ir arm\'andolo sin confundirse con los par\'ametros. Yo no lo implement\'e como
patr\'on formal, pero la misma idea est\'a presente en el formulario de
creaci\'on de combates en Flutter: el usuario llena campo por campo antes de
mandar todo al backend, y en el backend, \texttt{CombateService} tambi\'en
construye el resultado del combate de forma incremental.
 
\subsubsection*{Implementaci\'on en el Backend (Java / Spring Boot)}
 
El m\'etodo \texttt{simularCombate} es el mejor ejemplo: comienza con un
\texttt{Combate} base y lo va enriqueciendo paso a paso (c\'alculo de stats,
bonus por intensidad, bonus por tipo, bonus por lugar, evento aleatorio,
ganador final) antes de persistirlo.
 
\begin{lstlisting}[style=javaStyle, caption={Construcci\'on incremental del resultado en CombateService (service/CombateService.java)}]
public Combate simularCombate(Long id) {
    Optional<Combate> opt = repo.findById(id);
    if (opt.isEmpty()) return null;
    Combate c = opt.get();
 
    // Paso 1: stats base
    int ataque    = random.nextInt(100) + 1;
    int defensa   = random.nextInt(100) + 1;
    int velocidad = random.nextInt(100) + 1;
 
    // Paso 2: bonus segun intensidad
    switch (c.getIntensidad().toLowerCase()) {
        case "epico"  -> { ataque += 20; velocidad += 15; }
        case "alto"   -> { ataque += 10; velocidad += 5;  }
        case "medio"  ->   ataque += 5;
    }
 
    // Paso 3: bonus segun tipo de combate
    switch (c.getTipo().toLowerCase()) {
        case "supervivencia" -> defensa += 15;
        case "equipo"        -> { ataque += 10; defensa += 10; }
    }
 
    // Paso 4: bonus segun lugar
    switch (c.getLugar().toLowerCase()) {
        case "namek"  -> ataque    += 10;
        case "torneo" -> velocidad += 10;
        case "tierra" -> defensa   += 5;
    }
 
    // Paso 5: evento y ganador
    String evento  = generarEvento();
    String ganador = (ataque + velocidad) >= (defensa + random.nextInt(80))
            ? "Guerrero A" : "Guerrero B";
 
    // Paso 6: ensamblar el resultado final y persistir
    c.setResultado(
        "Lugar: "      + c.getLugar()      +
        " | Intensidad: " + c.getIntensidad() +
        " | Evento: "  + evento            +
        " | ATQ:" + ataque + " DEF:" + defensa + " VEL:" + velocidad +
        " | Ganador: " + ganador
    );
    return repo.save(c);
}
\end{lstlisting}
 
\subsubsection*{Reflejo en el Frontend (Dart / Flutter)}
 
En la pantalla \texttt{CrearCombateScreen} el usuario va llenando campo por
campo (a\~no, intensidad, tipo, lugar, creador, imagen), y solo cuando todo
est\'a listo se ensambla el objeto y se env\'ia:
 
\begin{lstlisting}[style=dartStyle, caption={Construcci\'on incremental del Combate en el formulario (screens/crear\_combate\_screen.dart -- fragmento)}]
// Cada campo se guarda en su variable de estado
String _lugar      = '';
String _intensidad = 'Medio';
String _tipo       = 'Individual';
int    _anio       = 2024;
String _creador    = '';
 
// Al confirmar, se ensambla el objeto completo
void _guardar() async {
  final combate = Combate(
    anio:       _anio,
    intensidad: _intensidad,
    tipo:       _tipo,
    lugar:      _lugar,
    creador:    _creador,
  );
  await _repo.crearCombate(combate);  // Se envia solo cuando esta completo
}
\end{lstlisting}
 
% ============================================================================
\section{Patrones Estructurales}
% ============================================================================
 
Estos patrones se enfocan en c\'omo se relacionan y se organizan las clases
entre s\'i. Sirven para que el c\'odigo sea m\'as flexible y f\'acil de mantener
sin tener que cambiar todo cada vez que algo cambia.
 
% ----------------------------------------------------------------------------
\subsection{Repository}
% ----------------------------------------------------------------------------
 
Este fue el que m\'as us\'e en el proyecto. La idea es poner una capa entre
la l\'ogica de la aplicaci\'on y el origen de los datos (que puede ser una
base de datos, una API, lo que sea).
 
En Flutter tengo una clase \texttt{CombateRepository} que es la \'unica
que sabe c\'omo hablar con el backend. Las pantallas no llaman directamente
al servicio HTTP, sino que le piden al repositorio. As\'i, si en alg\'un momento
cambia la URL o la forma de obtener datos, solo tengo que modificar el
repositorio y no tocar todas las pantallas.
 
\subsubsection*{Implementaci\'on en el Frontend (Dart / Flutter)}
 
\begin{lstlisting}[style=dartStyle, caption={CombateRepository -- capa de acceso a datos (lib/repositories/combate\_repository.dart)}]
import '../models/combate.dart';
import '../services/api_service.dart';
 
class CombateRepository {
  final _api = ApiService();   // Dependencia interna oculta al caller
 
  // Las pantallas solo conocen estos metodos, no saben nada de HTTP
  Future<List<Combate>> obtenerCombates()       => _api.obtenerCombates();
  Future<List<Combate>> obtenerMisCombates()    => _api.obtenerMisCombates();
  Future<Combate>       obtenerPorId(int id)    => _api.obtenerPorId(id);
  Future<Combate>       crearCombate(Combate c) => _api.crearCombate(c);
  Future<Combate>       actualizarCombate(int id, Combate c)
                                                => _api.actualizarCombate(id, c);
  Future<Combate>       actualizarImagen(int id, String url)
                                                => _api.actualizarImagen(id, url);
  Future<void>          eliminarCombate(int id) => _api.eliminarCombate(id);
  Future<Combate>       simularCombate(int id)  => _api.simularCombate(id);
 
  Future<Map<String, int>>     obtenerReacciones(int id)
                                                => _api.obtenerReacciones(id);
  Future<List<String>>         obtenerMisReacciones(int id)
                                                => _api.obtenerMisReacciones(id);
  Future<Map<String, dynamic>> toggleReaccion(int id, String t)
                                                => _api.toggleReaccion(id, t);
}
\end{lstlisting}
 
\subsubsection*{Implementaci\'on en el Backend (Java / Spring Boot)}
 
En el backend, Spring Data JPA aplica el mismo patr\'on de forma autom\'atica.
\texttt{CombateRepository} extiende \texttt{JpaRepository} y define consultas
personalizadas sin que el servicio sepa nada de SQL:
 
\begin{lstlisting}[style=javaStyle, caption={CombateRepository en Spring Data JPA (repository/CombateRepository.java)}]
package com.example.dragon.repository;
 
import com.example.dragon.entity.Combate;
import org.springframework.data.jpa.repository.JpaRepository;
import java.util.List;
 
public interface CombateRepository extends JpaRepository<Combate, Long> {
    // Spring genera automaticamente el SQL; el servicio solo llama al metodo
    List<Combate> findByUsuarioId(Long usuarioId);
}
\end{lstlisting}
 
% ----------------------------------------------------------------------------
\subsection{Adapter}
% ----------------------------------------------------------------------------
 
El Adapter sirve cuando tienes dos partes del sistema que no hablan el
mismo idioma y necesitas un intermediario que traduzca.
 
En mi proyecto lo vi reflejado en los m\'etodos \texttt{fromJson()} de los
modelos en Dart: el backend devuelve JSON con ciertos nombres de campos, y el
modelo se encarga de traducirlos al formato que usa Flutter internamente. De
forma equivalente, \texttt{toJson()} hace la traducci\'on en sentido contrario.
 
\subsubsection*{Implementaci\'on en el Frontend (Dart / Flutter)}
 
\begin{lstlisting}[style=dartStyle, caption={Adapter JSON$\leftrightarrow$Dart en el modelo Combate (lib/models/combate.dart)}]
class Combate {
  // ...campos internos de Flutter...
 
  // Adapter: convierte JSON del backend -> objeto Dart
  factory Combate.fromJson(Map<String, dynamic> j) => Combate(
    id:           j['id'],
    anio:         j['anio'],
    intensidad:   j['intensidad'],
    tipo:         j['tipo'],
    lugar:        j['lugar'],
    creador:      j['creador'],
    usuarioId:    j['usuarioId'],    // clave snake_case del backend
    resultado:    j['resultado'],
    fechaCreacion: j['fechaCreacion'],
    imagen: j['imagen'] != null
        ? ImagenCombate.fromJson(j['imagen'])
        : null,
  );
 
  // Adapter inverso: convierte objeto Dart -> JSON para el backend
  Map<String, dynamic> toJson() => {
    'anio':       anio,
    'intensidad': intensidad,
    'tipo':       tipo,
    'lugar':      lugar,
    'creador':    creador,
  };
}
\end{lstlisting}
 
\subsubsection*{Implementaci\'on en el Backend (Java / Spring Boot)}
 
En el backend, los DTOs de \texttt{AuthDtos} cumplen el mismo rol: adaptan
los datos que llegan del cliente al formato que usa la l\'ogica interna, y
viceversa para las respuestas.
 
\begin{lstlisting}[style=javaStyle, caption={Adapter entre capas con DTOs (dto/AuthDtos.java -- fragmento)}]
public class AuthDtos {
 
    // Adapter de entrada: traduce el JSON del cliente a un objeto Java
    public static class LoginRequest {
        private String email;
        private String password;
        public String getEmail()    { return email; }
        public String getPassword() { return password; }
    }
 
    // Adapter de salida: traduce el objeto Java al JSON que espera Flutter
    public static class AuthResponse {
        private final String token;
        private final Long   usuarioId;   // nombre que usa Flutter
        private final String nombre;
        private final String email;
 
        public AuthResponse(String token, Long usuarioId,
                            String nombre, String email) {
            this.token     = token;
            this.usuarioId = usuarioId;
            this.nombre    = nombre;
            this.email     = email;
        }
        // getters...
    }
}
\end{lstlisting}
 
% ----------------------------------------------------------------------------
\subsection{Facade}
% ----------------------------------------------------------------------------
 
El Facade es b\'asicamente crear una clase sencilla que oculte la complejidad
interna. Desde afuera ves una interfaz simple, pero por dentro puede estar
haciendo muchas cosas.
 
Mi \texttt{ApiService} del proyecto funciona exactamente as\'i: las pantallas
solo llaman m\'etodos como \texttt{obtenerCombates()} sin preocuparse por los
detalles HTTP, el manejo del token, la serializaci\'on JSON, etc. Todo eso
est\'a escondido dentro del \texttt{ApiService}.
 
\subsubsection*{Implementaci\'on en el Frontend (Dart / Flutter)}
 
\begin{lstlisting}[style=dartStyle, caption={ApiService -- Facade sobre HTTP (lib/services/api\_service.dart)}]
class ApiService {
  static const String baseUrl = 'https://dragonw.onrender.com';
  static const String _path   = '$baseUrl/api/combates';
  static const String _rpath  = '$baseUrl/api/reacciones';
 
  // Detalle interno: headers publicos
  static const Map<String, String> _headersPublic = {
    'Content-Type': 'application/json',
  };
 
  // Detalle interno: headers con JWT inyectado automaticamente
  Map<String, String> get _headers {
    final token = AuthService.instance.token;
    return {
      'Content-Type': 'application/json',
      if (token != null) 'Authorization': 'Bearer $token',
    };
  }
 
  // Interfaz simple que usan las pantallas: un solo metodo, cero detalles HTTP
  Future<List<Combate>> obtenerCombates() async {
    final res = await http.get(Uri.parse(_path));
    if (res.statusCode == 200) {
      final list = jsonDecode(res.body) as List;
      return list.map((e) => Combate.fromJson(e)).toList();
    }
    throw Exception('Error ${res.statusCode} al obtener combates');
  }
 
  // Toggle de reaccion: internamente hace POST, maneja token, decodifica JSON
  Future<Map<String, dynamic>> toggleReaccion(
      int combateId, String tipo) async {
    final res = await http.post(
      Uri.parse('$_rpath/$combateId/toggle'),
      headers: _headers,                       // JWT oculto al caller
      body: jsonEncode({'tipo': tipo}),
    );
    if (res.statusCode == 200) {
      final data   = jsonDecode(res.body) as Map<String, dynamic>;
      final conteo = data['conteo'] as Map<String, dynamic>;
      return {
        'added': data['added'] as bool,
        'tipo':  data['tipo']  as String,
        'like':  (conteo['like']  as num).toInt(),
        'heart': (conteo['heart'] as num).toInt(),
        'laugh': (conteo['laugh'] as num).toInt(),
      };
    }
    throw Exception('Error ${res.statusCode}: ${res.body}');
  }
  // ... todos los demas metodos siguen el mismo patron
}
\end{lstlisting}
 
Gracias al Facade, el c\'odigo de las pantallas queda limpio y enfocado solo
en mostrar informaci\'on:
 
\begin{lstlisting}[style=dartStyle, caption={Uso del Facade desde una pantalla (fragmento de home\_screen.dart)}]
// La pantalla no sabe nada de HTTP, tokens ni JSON
final combates = await _repo.obtenerCombates();
setState(() => _combates = combates);
\end{lstlisting}
 
% ============================================================================
\section{Conclusi\'on}
% ============================================================================
 
Al principio me parec\'ia que los patrones de dise\~no son muy dif\'iciles de
aplicar en proyectos reales, pero conforme fui entendiendo y practicando en
clase, y al mismo tiempo aplic\'andolos en mi proyecto, me di cuenta de que los
estaba usando sin saberlo, o que muchas cosas que hac\'ia por intuici\'on ten\'ian
un nombre y una raz\'on de ser.
 
No creo que haya que memorizar todos los patrones ni intentar usarlos todos
en un mismo proyecto. La clave es entender el problema que resuelve cada uno
y tenerlos en mente para cuando uno los necesite.
 
\end{document}

